\subsection{The remaining unrestricted problem}
All cases with Frobenius--Perron character field of degree at most two
are now classified. The proper-subring theorem also handles arbitrary
character degree when a proper nontrivial fusion subring is present.
The known rings covered by these two results are the fourteen entries of
the realization table other than $F_{1,6}$, the ring of
$\operatorname{PSU}(2)_7$.

\begin{question}\label{question:unrestricted-simple}
Let $R$ be a self-dual rank-four fusion ring with no proper nontrivial
fusion subring and with Frobenius--Perron character field of degree three
or four. If $R$ admits a pivotal categorification, must it be the ring of
$\operatorname{PSU}(2)_7$?
\end{question}

No complete answer is obtained here. A pivotal realization of a ring in
this question has a pseudo-unitary Galois conjugate, by
Corollary~\ref{cor:simple-pivotal-galois-pu}. Cyclotomicity makes its
character field cyclic cubic, cyclic quartic or biquadratic. In the cubic
case, Proposition~\ref{prop:cubic-rational-codegree-bound} further restricts
the rational character to formal codegree three or four. The finite
classification constrains any additional positive ring to lie outside
its certified range, but supplies no upper bound for these remaining
field types.

Section~\ref{sec:cubic-codegree-three-difference} gives a
further uniform reduction when the rational character in the cubic case
has codegree three. At most six central simples distinguish two specified
basis objects, and one possible pattern is excluded by Galois symmetry.
Their dimensions remain unbounded. In particular, a bound on the number
of terms in a signed twist identity is not a bound on its cyclotomic
length. These constraints do not settle the remaining question, and no
unrestricted classification of all rank-four pivotal rings is claimed.
